\documentclass[12pt]{article}

\usepackage{amsmath}
\usepackage{geometry}
\geometry{margin=1in}

\setlength{\parskip}{0.5em}

\begin{document}

\title{Kinematics: Projection Practice 2}
\author{}
\date{}
\maketitle

\vspace{1em}

\textbf{Instructions:}

\begin{enumerate}
\item Describe the motion before writing formulas.
\item Show all steps clearly.
\item Include units in every step and in the final answer.
\end{enumerate}

\noindent\rule{\textwidth}{0.4pt}

\newpage

\section*{Projectile Motion Practice Problems (Set 2)}

\subsection*{1. Angled Launch on Level Ground}
A soccer ball is kicked from the ground with an initial velocity of $25 \text{ m/s}$ at an angle of $37^\circ$ above the horizontal. Assume the ground is level and neglect air resistance. (Take $g = 10 \text{ m/s}^2$, $\sin 37^\circ = 0.6$, $\cos 37^\circ = 0.8$) \\

\textbf{Find:} The total time the ball is in the air, the maximum height it reaches, and its maximum horizontal range.

\newpage

\subsection*{2. Launch from an Elevated Position}
A stone is thrown with an initial speed of $15 \text{ m/s}$ at an angle of $30^\circ$ above the horizontal from the edge of a cliff that is $45 \text{ m}$ high. (Take $g = 10 \text{ m/s}^2$, $\sin 30^\circ = 0.5$, $\cos 30^\circ \approx 0.866$) \\

\textbf{Find:} The time it takes for the stone to reach the ground below and the horizontal distance from the base of the cliff to where the stone lands.

\newpage

\subsection*{3. Clearing an Obstacle}
An archer shoots an arrow with an initial velocity of $40 \text{ m/s}$ at an angle of $45^\circ$ to the horizontal. There is a $15 \text{ m}$ tall vertical wall located $120 \text{ m}$ away from the archer. (Take $g = 10 \text{ m/s}^2$, $\sin 45^\circ = \cos 45^\circ = \frac{\sqrt{2}}{2} \approx 0.707$) \\

\textbf{Find:} Will the arrow pass over the wall? Support your answer with calculations of the arrow's height when it reaches the horizontal position of the wall.

\newpage

\subsection*{4. Conceptual Projectile Motion}
For a projectile launched at an angle $\theta$ ($0^\circ < \theta < 90^\circ$) from level ground, fill in the blanks: \\
The horizontal component of its velocity remains \underline{\hspace{3cm}} throughout the flight. As the projectile moves towards its peak, the magnitude of its vertical velocity \underline{\hspace{3cm}}. At the highest point of the trajectory, its vertical velocity is \underline{\hspace{2cm}} $\text{ m/s}$, and its acceleration is \underline{\hspace{3cm}} directed \underline{\hspace{3cm}}. 

\vspace{10cm}

\subsection*{5. Maximum Height vs. Range}
A projectile is launched from ground level with an initial velocity $v_0$ at a launch angle $\theta$. It is observed that the projectile's maximum height is exactly equal to its total horizontal range. \\

\textbf{Find:} The launch angle $\theta$. (You may express your answer using an inverse trigonometric function, e.g., $\theta = \arctan(x)$).

\end{document}
